% ==========================================================
% 相同球放入有标号盒，且数量互不相同
% 难度：★★ 中档
% ==========================================================

\newcommand{\qProbPermDiffBoxesStem}{%
将 $8$ 个完全相同的球放入编号为 $1,2,3$ 的三个空盒中，要求放入后三个盒子均不空，且三个盒子中的球数均不相同。求放法数。
}

\newcommand{\qProbPermDiffBoxesAnswer}{%
$12$
}

\newcommand{\qProbPermDiffBoxesSolution}{%
设三个盒中的球数分别为 $x_1,x_2,x_3$。则 $x_1+x_2+x_3=8$，并且 $x_1,x_2,x_3$ 是互不相同的正整数。

由于球完全相同，一种放法只取决于各盒的球数。先暂时忽略盒子编号，把三个数从小到大排列：$a<b<c$ 且 $a+b+c=8$。

从最小数开始枚举：
\begin{itemize}
    \item 当 $a=1$ 时，可以得到 $(1,2,5),\ (1,3,4)$；
    \item 当 $a\geqslant2$ 时，最小可能和为 $2+3+4=9>8$，不再可能。
\end{itemize}

因此无序数量类型共有两种。每种数量类型中的三个数互不相同，可以任意分配给编号为 $1,2,3$ 的三个盒子，各有 $3!$ 种分配。

故总放法数为 $2\times3!=12$。
}

\newcommand{\qProbPermDiffBoxesTeachingNotes}{%
\textbf{【避坑与边界说明】}
\begin{itemize}
    \item 球完全相同，不能选择“哪几个球”进入某个盒子。
    \item 盒子有编号，因此同一组数量的不同排列是不同放法。
    \item 不能直接用 $\binom{7}{2}$，因为该公式还包括两个或三个盒子中数量相同的情况。
    \item 本题每种数量类型中的三个数都不同，所以乘 $3!$；若有两个数相同，只能乘 $\frac{3!}{2!}$。
\end{itemize}
}
